\section{Methodology}

We designed a controlled study to isolate how the composition of inference-time Self-Refinement protocols affects code correctness and token consumption.
All protocols used a Shared-Initial Design, in which the same model revised the same Self-Generated Initial Candidate without access to tests, execution results, compiler diagnostics, runtime traces, or external verifier feedback.

\subsection{Research Questions}

\researchquestion{RQ1}{How does the stage composition of Self-Refinement affect final code correctness?}
Explicit Critique Generation and Revision Planning may structure a model's review, but each additional stage is useful only if it produces a better Final Candidate.
We isolated their incremental effects through paired comparisons of Direct Generation with Direct Revision (\texttt{R}), \texttt{R} with Critique-Conditioned Revision (\texttt{CR}), and \texttt{CR} with Critique-and-Plan-Conditioned Revision (\texttt{CPR}) on the same Initial Candidates.

\researchquestion{RQ2}{How do alternative refinement protocols balance repairs of initially incorrect candidates and regressions of initially correct candidates?}
Final Pass Rate alone can conceal whether a protocol repairs many failures while also damaging correct programs.
We therefore classified every evaluable initial-to-final pair as Repair, Correct-to-Incorrect Regression, Functional Preservation, or Unrepaired Failure and compared these paired transitions across protocols.

\researchquestion{RQ3}{How does the value of Decision-Conditioned Refinement vary across model-benchmark combinations with different Initial Pass Rates?}
A preservation decision has different consequences when most Initial Candidates are already correct than when many remain repairable, creating a tradeoff between preventing regressions and missing repairs.
We applied one independently generated Refinement-Need Decision to each Always-Refine protocol and related the resulting correctness and token-cost differences to the Initial Pass Rate of each of the 18 model-benchmark combinations.

\researchquestion{RQ4}{Do the correctness effects of additional refinement stages justify their token costs?}
Critique, planning, and decision calls can improve a candidate or avoid an unnecessary revision, but they also add input and output tokens.
We compared Protocol Token Cost with final correctness at both candidate and benchmark levels, including avoided refinement tokens, net token savings, and correctness-cost Pareto relationships within each model.

\subsection{Self-Refinement Protocols}

For each model-task pair, Direct Generation produced one Initial Candidate from the Task Specification.
The candidate was stored once and reused as the exact input to all subsequent Refinement Stages, so protocol contrasts were not confounded by different initial generations.
Benchmark tests were reserved for post-hoc evaluation and were never included in a model prompt.

The protocols were composed from four observable stages, each implemented as a separate prompt and model call:

\begin{enumerate}
    \item \textbf{Refinement-Need Decision (D).} Given the Task Specification and Initial Candidate, the model returned \texttt{PRESERVE} when the candidate should be retained exactly or \texttt{REFINE} when functional revision was needed.
    \item \textbf{Critique Generation (C).} The model identified potential functional mismatches, missing conditions, and relevant fault locations, or recommended preservation when it found no problem.
    \item \textbf{Revision Planning (P).} Given the same candidate and its stored critique artifact, the model converted the critique into an actionable set of changes while identifying behavior that should remain intact.
    \item \textbf{Code Revision (R).} The model returned a complete Revised Candidate using the Task Specification, Initial Candidate, and the intermediate artifacts available under the protocol.
\end{enumerate}

Critique Generation and Revision Planning were instructed not to provide a complete revised program, which kept diagnosis, planning, and code production as distinct calls.
Conversely, every Code Revision prompt permitted the model to retain the Initial Candidate when no functional change was needed.

Table~\ref{tab:protocols} summarizes the Stage Composition of the baseline and three Always-Refine protocols.
\texttt{R} revises the candidate directly, \texttt{CR} adds an explicit critique before revision, and \texttt{CPR} further adds a revision plan between critique and revision.

\begin{table}[H]
\vspace*{1.5cm}
\caption{Protocol composition and the information supplied to the final Code Revision call.}
\label{tab:protocols}
\begin{tabularx}{\textwidth}{p{0.18\textwidth}p{0.23\textwidth}L}
\toprule
Protocol & Composition after Direct Generation & Final-candidate construction \\
\midrule
Direct Generation & None & Use the Initial Candidate without refinement. \\
\texttt{R} & $R$ & Revise from the Task Specification and Initial Candidate. \\
\texttt{CR} & $C \rightarrow R$ & Revise from the Initial Candidate and stored critique artifact. \\
\texttt{CPR} & $C \rightarrow P \rightarrow R$ & Revise from the Initial Candidate, stored critique artifact, and revision plan. \\
\bottomrule
\end{tabularx}
\end{table}

Decision-Conditioned Refinement combined the Decision with each corresponding Always-Refine result.
If the Decision was \texttt{PRESERVE}, the Final Candidate was the exact stored Initial Candidate; if it was \texttt{REFINE}, the Final Candidate was the Revised Candidate from \texttt{R}, \texttt{CR}, or \texttt{CPR}.
This construction yielded Decision-Conditioned Direct Revision (\texttt{DR}), Decision-Conditioned Critique-Conditioned Revision (\texttt{DCR}), and Decision-Conditioned Critique-and-Plan-Conditioned Revision (\texttt{DCPR}) without additional revision calls.
The Decision was not shown the protocol name, critique, plan, or any execution evidence, and its output was not passed into the refinement prompts.

\subsection{Models and Benchmarks}

We selected publicly available open-weight checkpoints that could be served under one two-GPU workstation budget while covering different model families, sizes, and dense and mixture-of-experts architectures.
The panel emphasized code-specialized instruction models and included Gemma 4 as a general-purpose control; all within-model protocol comparisons used an unchanged checkpoint and precision.
Table~\ref{tab:models} lists the exact checkpoints and their roles.

\begin{table}[htbp]
\vspace*{1.5cm}
\caption{Open-weight models used in the study.}
\label{tab:models}
\begin{tabularx}{\textwidth}{p{0.29\textwidth}p{0.22\textwidth}p{0.12\textwidth}L}
\toprule
Model and checkpoint & Architecture and size & Precision & Role \\
\midrule
Qwen2.5-Coder-7B\newline{\scriptsize\nolinkurl{Qwen/Qwen2.5-Coder-7B-Instruct}} & Code-specialized dense, approximately 7.6B & BF16 & Smaller code-specialized model. \\
Qwen2.5-Coder-14B\newline{\scriptsize\nolinkurl{Qwen/Qwen2.5-Coder-14B-Instruct}} & Code-specialized dense, 14.7B & BF16 & Within-family scale comparison. \\
DeepSeek-Coder-V2-Lite\newline{\scriptsize\nolinkurl{deepseek-ai/DeepSeek-Coder-V2-Lite-Instruct}} & Code-specialized MoE, 16B total and 2.4B active & BF16 & Efficient MoE model from a different family. \\
Devstral-Small-2-24B\newline{\scriptsize\nolinkurl{mistralai/Devstral-Small-2-24B-Instruct-2512}} & Software-engineering-specialized dense, 24B & Official FP8 & Larger dense code model. \\
Qwen3-Coder-30B-A3B\newline{\scriptsize\nolinkurl{Qwen/Qwen3-Coder-30B-A3B-Instruct-FP8}} & Code-specialized MoE, 30.5B total and 3.3B active & Official FP8 & Larger code-specialized MoE model. \\
Gemma-4-31B-QAT-W4A16\newline{\scriptsize\nolinkurl{google/gemma-4-31B-it-qat-w4a16-ct}} & General-purpose dense, 31B & Official QAT W4A16 & General-purpose instruct control. \\
\bottomrule
\end{tabularx}
\end{table}

The differing precisions reflect the official locally deployable checkpoints that fit the common hardware budget, rather than a protocol-level treatment.
We therefore interpret raw score and token differences primarily within a model, not as precision-independent rankings across models or as equivalent compute across tokenizers.
A planned StarCoder2-15B-Instruct model was removed before the primary experiment after three pilots showed repeated failures to follow the candidate-fence and categorical-Decision contracts.

We chose three Python code-generation benchmarks to vary specification style, library use, and task composition without assigning them a priori difficulty labels.
Except for exclusions fixed before model inference for evaluator compatibility or offline reproducibility, we used every task in the pinned releases, as summarized in Table~\ref{tab:benchmarks}.

\begin{table}[htbp]
\caption{Benchmarks and task scope.}
\label{tab:benchmarks}
\begin{tabularx}{\textwidth}{p{0.23\textwidth}p{0.12\textwidth}p{0.27\textwidth}L}
\toprule
Benchmark & Tasks used & Task format & Evaluation \\
\midrule
HumanEval+ v0.1.10 & 163 of 164 & Function signatures and docstrings & EvalPlus 0.3.1 augmented tests. \\
MBPP+ v0.2.0 & 378 of 378 & Natural-language problems with examples & EvalPlus 0.3.1 augmented tests. \\
BigCodeBench-Instruct v0.1.4 & 1,136 of 1,140 & Natural-language instructions with diverse library and compositional requirements & BigCodeBench 0.2.5 unit tests. \\
\bottomrule
\end{tabularx}
\end{table}

We excluded \texttt{HumanEval/32} because its pinned canonical solution was incompatible with the EvalPlus special oracle and success bookkeeping.
We also excluded \texttt{BigCodeBench/101}, \texttt{/590}, \texttt{/1005}, and \texttt{/1012} because they required live external content that could not be reproduced with outcome-neutral offline resources.
The resulting scope contained 1,677 tasks per model and 10,062 model-task pairs across the six models.
The Initial Pass Rate of each model-benchmark combination was treated as an observed continuous characteristic rather than using fixed easy, medium, and hard benchmark labels.

\subsection{Experimental Setting}

\paragraph{Phase-wise inference and artifact reuse.}
We executed the experiment by model-resident phase rather than running one task through every protocol end to end.
For each checkpoint, the model remained loaded while we processed all 1,677 tasks in the following order: Direct Generation, Decision, Direct Revision, Critique Generation, Critique-Conditioned Revision, Revision Planning, and Plan-Conditioned Revision.
Phase boundaries checkpointed immutable artifacts, allowing completed logical calls to be resumed without regeneration.
The same critique artifact was reused by \texttt{CR} and Revision Planning, the same plan was reused by \texttt{CPR}, and one Decision was reused to derive \texttt{DR}, \texttt{DCR}, and \texttt{DCPR}.
This schedule reduced model-loading and duplicate-call costs while preserving the dependencies and inputs of the conceptual protocols.

\paragraph{Inference environment.}
Inference ran on an Ubuntu 22.04 server with two NVIDIA GeForce RTX 4090 GPUs, each with 24~GiB of memory.
We used vLLM 0.26.0 with tensor parallelism of two, a batch size of 32, a 16,384-token context limit, and GPU memory utilization set to 0.85.
All prompts used the checkpoint's tokenizer and chat template as direct token-ID inputs, automatic prompt truncation was disabled, and the rendered prompt plus output allowance was required to fit the common context limit.
Generation used the PyTorch-native sampler with batch-invariant execution, one output, greedy decoding (temperature 0, top-$p$ 1, top-$k$ 0), seed 0, and no frequency, presence, or repetition penalties.
Maximum output lengths were 4,096 tokens for Direct Generation and revision calls, 2,048 for Critique Generation and Revision Planning, and 64 for the Decision.

\paragraph{Output contracts and observed exceptions.}
Candidate prompts required exactly one complete fenced Python block; text outside that block was ignored, whereas responses with no complete block, multiple blocks, or an incomplete block were not repaired or heuristically extracted.
Across 70,386 executed model calls, eight Direct Generation responses and ten revision responses violated this contract and were retained as \texttt{malformed\_candidate}; the eight malformed Initial Candidates blocked their 48 dependent calls.
These cases remained in the study scope and counted as zero only for the end-to-end success measure, while functional transitions were left unevaluated because no candidate was executed.
For Decisions, the primary parser accepted only a complete response equal to \texttt{PRESERVE} or \texttt{REFINE} after trimming surrounding whitespace.
Of 202 non-exact responses, 156 differed only by one terminal period and were normalized deterministically, while 46 directionally unambiguous explanations were coded before candidate evaluation under a fixed, execution-blind rubric; all were resolved, and the original invalid responses and token counts were retained.
We additionally analyze exact-parsed Decisions alone to assess sensitivity to this resolution procedure.

\paragraph{Candidate evaluation.}
Initial and Final Candidates were evaluated in processes separated from model inference, and a candidate passed only if it satisfied the benchmark's full test suite.
The orchestration, analysis, and EvalPlus environments used Python 3.12, whereas BigCodeBench execution used a separately pinned Python 3.10.12 environment for compatibility with its official dependency range.
Candidate execution was isolated from the host process namespace and external network; pinned local resources and narrowly scoped offline fixtures supported reproducible BigCodeBench tasks without exposing execution results to later model calls.
EvalPlus candidates were additionally subject to an 8~GiB process-memory limit.

Before model inference, we ran each official reference solution three times in its final evaluator environment to fix task-level timeouts independently of generated outputs.
The default primary wall-time limit was 10 seconds; the audit assigned a 60-second limit to \texttt{Mbpp/599} and limits from 30 to 240 seconds to 15 BigCodeBench tasks, with no HumanEval+ primary overrides.
After the complete primary evaluation batch, a provisional timeout was retried exactly once using the larger of 120 seconds and 1.5 times its primary limit; \texttt{HumanEval/38}, \texttt{/50}, and \texttt{/53} used a pre-specified 180-second confirmation limit to accommodate EvalPlus's internal bounds.
We assigned final \texttt{TIMEOUT} only when both attempts timed out, and an independently observed assertion or exception took precedence as functional \texttt{FAIL}.

\paragraph{Outcome and cost analysis.}
We computed Initial and Final Pass Rates and decomposed each evaluable paired outcome into Repair, Correct-to-Incorrect Regression, Functional Preservation, or Unrepaired Failure.
Protocol comparisons used a task-level paired bootstrap with 10,000 resamples, a fixed seed of 20260804, and 95\% confidence intervals.
Token accounting retained input and output tokens for every call and summed only the stages required by the protocol: \texttt{R} used its revision call, \texttt{CR} used critique plus critique-conditioned revision, and \texttt{CPR} used critique, plan, and plan-conditioned revision.
For a Decision-Conditioned protocol, cost consisted of the Decision alone after \texttt{PRESERVE}, or the Decision plus the corresponding Always-Refine cost after \texttt{REFINE}.
